\section{Conclusion}
\label{sec:conclusion}

We have proposed that two large, parallel online discourses---``AI overuses em dashes'' and ``AI thinks in markdown''---describe the same phenomenon at different levels of compression. The em dash is not an arbitrary stylistic tic. It is the smallest surviving unit of the structural orientation that LLMs acquire from markdown-saturated training corpora: a formatting artifact that persists under output constraints because it is simultaneously valid punctuation and structural architecture.

The five-step genealogy---training data saturation, structural internalization, the dual-register status of the em dash, post-training amplification, and the complementarity of training and fine-tuning explanations---provides an explanatory account that unifies these observations. Empirical measurements across twelve models from five providers confirm the predicted suppression asymmetry: markdown features are eliminated or nearly eliminated by prose instructions, while em dashes persist at rates determined by the specific fine-tuning procedure applied. A three-condition suppression gradient demonstrates the dual-register mechanism directly: generic formatting suppression does not catch the em dash, and even explicit prohibition fails to eliminate it in some models. A base-vs-instruct comparison confirms that a latent tendency exists pre-RLHF, with the direction and magnitude of amplification determined by the fine-tuning procedure.

The em dash functions not merely as a tell of AI-generated text, but as a diagnostic signature of how a model was fine-tuned---varying from zero (Llama) to near-invariant under explicit prohibition (GPT-4.1 at 6.97 per 1,000 words in long-form output). Once named, this mechanism becomes straightforward to address, as the variation across providers and model generations already demonstrates.
